\chapter{Biết Đủ Là Giàu Bên Trong}
\markboth{Biết Đủ Là Giàu Bên Trong}{Knowing Enough Is Inner Wealth}

\section{Bao Giờ Mới Thấy Đủ?}
\begin{truyen}{Một Câu Hỏi Làm Người Ta Im Lặng}{A Question That Makes a Person Fall Silent}
\chuhoa{M}{ột người bạn hỏi người thầy rằng bao giờ thì anh mới thấy đủ. Câu hỏi đơn giản nhưng khiến anh ngồi lặng rất lâu.}
\textit{(A friend asked the teacher when he would finally feel that he had enough. The question was simple, but it left him silent for a long time.)}

Đời sống luôn có khả năng bày ra một thiếu hụt mới. Khi chưa có xe, người ta muốn có xe. Có xe rồi lại muốn xe tốt hơn. Khi chưa có nhà rộng, người ta muốn rộng hơn. Khi đã khá hơn một chút, lòng lại chỉ vào những thứ mình chưa với tới. Nếu cứ sống theo chiều ấy, cảm giác thiếu sẽ không bao giờ chấm dứt.
\textit{(Life always has the ability to present a new lack. When one has no vehicle, one wants one. After getting it, one wants a better one. When the house is small, one wants it bigger. Once life improves a little, the heart points toward what is still out of reach. If one keeps living in that direction, the feeling of shortage never ends.)}

Người thầy chợt hiểu rằng biết đủ không có nghĩa là thôi cố gắng. Nó có nghĩa là biết dừng lòng tham ở một giới hạn để sự biết ơn còn có chỗ mà đứng. Không có chỗ đứng cho biết ơn thì dù cầm nhiều tiền hơn trong tay, người ta vẫn nghèo ở bên trong.
\textit{(The teacher suddenly understood that knowing enough does not mean giving up effort. It means stopping greed at a certain boundary so that gratitude still has room to stand. Without room for gratitude, even more money in hand leaves a person inwardly poor.)}
\end{truyen}

\begin{baihoc}
Người không biết đủ thường rất khó thấy mình đang có gì, nên cũng rất khó giữ được những gì mình có.
\textit{(A person who does not know enough rarely sees what he already has, and therefore struggles to keep it.)}
\end{baihoc}

\begin{ghinhoanh}
Gratitude protects contentment.

Greed keeps the heart restless.
\end{ghinhoanh}

\begin{tuhocnhanh}
1. Điều gì em đang có nhưng lâu rồi không còn biết ơn? \textit{What do you already have that you have not appreciated for a long time?}

2. Thiếu thật sự của em là thiếu vật chất hay thiếu sự biết đủ? \textit{Is your real lack material, or is it a lack of contentment?}
\end{tuhocnhanh}
\ngancach

\section{Tri Túc Giả Phú}
\begin{truyen}{Người Biết Đủ Là Người Giàu}{He Who Knows Enough Is Rich}
\chuhoa{C}{ổ nhân nói ``tri túc giả phú'' --- người biết đủ là người giàu. Sự giàu này nghe qua tưởng như tinh thần, nhưng thật ra nó ảnh hưởng rất thật đến đời sống tiền bạc.}
\textit{(The ancients said, ``He who knows enough is rich.'' This kind of wealth sounds spiritual at first, but it has very real effects on financial life.)}

Người biết đủ không có nghĩa là không phấn đấu. Họ vẫn làm việc, vẫn cải thiện cuộc sống, vẫn nuôi ước mơ cho con cái. Nhưng khác ở chỗ lòng họ không bị thiếu liên tục xô đẩy. Vì thế, mỗi đồng tiền ở lại trong tay họ lâu hơn.
\textit{(A person who knows enough is not without ambition. Such a person still works, improves life, and dreams for the children. The difference is that his heart is not being constantly pushed by a sense of lack. Because of that, money stays with him longer.)}

Giàu bên trong là khi con người có thể nhìn vào căn nhà nhỏ của mình, vào bữa cơm giản dị của mình, vào đôi dép còn dùng tốt của mình và không xấu hổ. Sự không xấu hổ ấy cứu được rất nhiều đồng tiền khỏi những cuộc mua sắm vô nghĩa.
\textit{(To be inwardly rich is to look at one’s modest house, simple meal, and still-usable sandals without shame. That lack of shame saves a great deal of money from meaningless purchases.)}
\end{truyen}

\begin{baihoc}
Biết đủ không cản trở việc đi lên; nó ngăn việc đi lên bằng một trái tim luôn đói.
\textit{(Knowing enough does not block progress; it prevents progress from being driven by a perpetually hungry heart.)}
\end{baihoc}

\begin{ghinhoanh}
Contentment is not laziness.

It is freedom from endless hunger.
\end{ghinhoanh}

\begin{tuhocnhanh}
1. Em có nhầm lẫn giữa biết đủ và không cầu tiến không? \textit{Are you confusing contentment with lack of ambition?}

2. Điều gì giúp em thấy mình giàu hơn trong lòng? \textit{What helps you feel inwardly richer?}
\end{tuhocnhanh}
\ngancach

\section{Đừng Dùng Tiền Vá Một Khoảng Trống Tinh Thần}
\begin{truyen}{Khoảng Thiếu Không Món Đồ Nào Lấp Được}{A Gap No Purchase Can Fill}
\chuhoa{C}{ó những ngày người ta thấy lòng hụt hẫng, mệt mỏi, chán nản, rồi mở điện thoại ra mua một thứ gì đó cho khuây. Lúc bấm nút thanh toán, tưởng như mình vừa giải quyết xong nỗi nặng lòng.}
\textit{(There are days when a person feels empty, tired, and discouraged, then opens the phone and buys something for relief. At the moment of payment, it feels as though the heaviness has been solved.)}

Nhưng phần lớn những khoảng trống tinh thần không thể lấp bằng mua sắm. Chúng chỉ được che trong một lúc rất ngắn, rồi quay lại cùng với cảm giác tiếc tiền. Người thầy nhận ra càng cô đơn, càng mệt, càng dễ tiêu sai. Vì thế, có lúc giữ tiền không bắt đầu từ chiếc ví mà bắt đầu từ việc chăm sóc lòng mình cho lành hơn.
\textit{(But most inner emptiness cannot be filled by shopping. It is merely covered for a brief moment, then returns together with regret over the money spent. The teacher realized that the lonelier and more exhausted he felt, the easier it was to spend wrongly. In that sense, protecting money sometimes begins not with the wallet, but with caring for the heart more gently.)}

Biết đủ vì vậy cũng là biết nhận ra khi nào mình đang mua vì thiếu thốn vật chất, và khi nào mình đang mua vì thiếu một sự an ủi nào khác.
\textit{(Knowing enough therefore includes recognizing when we are buying because of material need, and when we are buying because we are missing some other form of comfort.)}
\end{truyen}

\begin{baihoc}
Không phải nỗi buồn nào cũng nên được chữa bằng mua sắm.
\textit{(Not every sadness should be treated with shopping.)}
\end{baihoc}

\begin{ghinhoanh}
Do not spend to soothe every ache.

Some emptiness needs another kind of care.
\end{ghinhoanh}

\begin{tuhocnhanh}
1. Em thường mua sắm nhiều hơn khi buồn hay khi mệt? \textit{Do you spend more when sad or when tired?}

2. Có khoảng trống nào trong lòng em đang bị mua sắm che tạm không? \textit{Is there an inner emptiness in you that shopping is only temporarily covering?}
\end{tuhocnhanh}

\begin{turanminh}
1. Tôi phải học biết đủ, nếu không dù kiếm thêm bao nhiêu tôi vẫn thấy thiếu và vẫn sống mệt.

2. Tôi không được dùng mua sắm để chữa mọi nỗi buồn, vì nhiều nỗi trống trong lòng chỉ bị che đi chứ không được chữa lành.

3. Biết đủ không phải là bỏ cố gắng. Đó là giữ cho lòng tham không cướp mất sự biết ơn và sự bình yên của mình.
\end{turanminh}

\begin{dayhoc}
1. Đây là chương rất phù hợp để dạy học sinh về lòng biết ơn và sự khác nhau giữa ước mơ lành mạnh với lòng tham không đáy.

2. Thầy có thể mời các em kể một điều mình đang có nhưng lâu rồi không còn trân trọng nữa.

3. Bài học nên chốt là: ``Người biết đủ không yếu chí; họ chỉ không để lòng mình đói mãi.''
\end{dayhoc}